\documentclass[11pt]{article}

\usepackage[margin=1.2in]{geometry}
\usepackage{amsmath, amssymb}
\usepackage{booktabs}
\usepackage{graphicx}
\usepackage{xcolor}
\usepackage{listings}
\usepackage{hyperref}

\newcommand{\R}{\mathbb{R}}
\newcommand{\Z}{\mathbb{Z}}
\newcommand{\Q}{\mathbb{Q}}
\newcommand{\T}{\mathbb{T}}
\newcommand{\Hb}{\mathbb{H}}
\newcommand{\re}{\operatorname{Re}}
\newcommand{\im}{\operatorname{Im}}

\definecolor{codebg}{rgb}{0.96, 0.96, 0.96}
\definecolor{codecomment}{rgb}{0.4, 0.5, 0.4}
\definecolor{codekw}{rgb}{0.2, 0.2, 0.6}

\lstset{
    language=C,
    basicstyle=\small\ttfamily,
    backgroundcolor=\color{codebg},
    commentstyle=\color{codecomment}\itshape,
    keywordstyle=\color{codekw}\bfseries,
    morekeywords={vec2,vec3,ivec3,float,int,bool,void,const,out},
    frame=single,
    framerule=0.4pt,
    xleftmargin=1em,
    xrightmargin=1em,
    aboveskip=1em,
    belowskip=1em,
    breaklines=true,
    columns=flexible,
}

\title{Raytracing Starscapes}
\author{}
\date{}

\begin{document}
\maketitle


%% ================================================================
\section{Making It a Shader Problem}
%% ================================================================

We've always rendered starscapes the same way: loop over polynomials, solve each one, plot the roots.
It works, but you're iterating over coefficient space and hoping the roots land somewhere useful.
Most of them fall off-screen, and the cost is something like $(2N+1)^{d+1}$, so it gets painful fast once you push the degree or coefficient bound up.

I want to describe a completely different approach that we've been playing with, where instead of looping over polynomials, we let each pixel ask ``am I near a root?'' on its own.


The idea is simple: for each pixel $z$ in the upper half plane, ask whether there's an integer polynomial of the right type with a root near $z$.
Black if yes, white if no.
That's a fragment shader --- one function call per pixel, all running in parallel on the GPU.

We only compute what's on screen, at whatever resolution we want, live.
The tradeoff is we lose marker shapes (the ``dots'' around each root are determined by closeness in coefficient space, so they're not round), and we rediscover the same polynomial from every nearby pixel.
But that redundancy costs nothing in a shader --- it's just parallel threads.


%% ================================================================
\section{Raytracing in the 3-Torus}
\label{sec:quadratic}
%% ================================================================

\subsection{The quadratic case}

Fix a pixel $z = s + iy$ in the upper half plane.
Every real quadratic $ax^2 + bx + c$ having $z$ as a root is a scalar multiple of
\[
    x^2 - 2sx + n, \qquad \text{where } s = \re(z),\; n = |z|^2.
\]
So the set of all such quadratics, parameterized by $t \in \R \setminus \{0\}$, is the ray
\[
    \mathbf{R}(t) = t \cdot (1,\; -2s,\; n)
\]
in the three-dimensional coefficient space $(a, b, c) \in \R^3$.
The pixel $z$ is close to a quadratic irrational if and only if this ray passes close to a point of the integer lattice $\Z^3$.

\subsection{A worked example}

Here's what the march actually does at $z = 0.01 + i$.
We have $s = 0.01$, $n = |z|^2 = 1.0001$, so the ray direction is
\[
    \mathbf{d} = (1,\; -0.02,\; 1.0001),
\]
close to but not exactly $(1, 0, 1)$, which would be the polynomial $x^2 + 1$.

Starting at $t_{\min} = 0.5$, we're at $(0.5,\; -0.01,\; 0.50)$, far from anything --- the SDF gives $\sim 0.67$, so we jump ahead by about $0.47$.
Now $t \approx 0.97$ and we're near $(0.97,\; -0.019,\; 0.97)$, whose nearest lattice point is $(1, 0, 1)$.
Distance is about $0.044$, just outside the ball radius of $0.04$, so we take a tiny step.
A few more sub-steps and we're inside the ball.

The lattice point is $(a, b, c) = (1, 0, 1)$, discriminant $0 - 4 = -4 < 0$, so it's a valid quadratic with complex roots --- the polynomial $x^2 + 1$, root at $i$.
Pixel colored black.

The closest approach to $(1, 0, 1)$ is at $t \approx 1.0$, where the offset is roughly $(0, -0.02, 0)$ --- the miss is almost entirely in the $b$-direction, magnitude $\approx 0.02$, well inside the $0.04$ ball.
A pixel farther from $i$ would miss by more than $0.04$, and the march would keep going.

\subsection{From rays to the 3-torus}

The question ``does the ray $t \cdot \mathbf{d}$ pass near an integer point?'' is the same as asking whether the fractional part $\{t \cdot \mathbf{d}\} = t \cdot \mathbf{d} \bmod \Z^3$ ever comes close to the origin in the 3-torus $\T^3 = \R^3 / \Z^3$.

In the torus picture, the ray wraps around and we're looking for returns to a neighborhood of the origin.
This is standard SDF raymarching: define
\[
    \operatorname{sdf}(\mathbf{p}) = |\mathbf{p} - \lfloor \mathbf{p} \rceil| - \varepsilon,
\]
where $\lfloor \cdot \rceil$ rounds to the nearest integer and $\varepsilon$ is the ball radius.
When $\operatorname{sdf} < 0$ we're inside a ball.
When $\operatorname{sdf} > 0$ it's a lower bound on the distance to the nearest lattice point, so we can safely step forward by $\operatorname{sdf} / |\mathbf{d}|$ --- largest safe step at every iteration.

\subsection{The shader}

Here's the quadratic shader:

\begin{lstlisting}
float sdfLatticeBalls(vec3 p, float radius) {
    return length(p - round(p)) - radius;
}

// Check that (a,b,c) gives a root in the upper half plane.
// disc < 0 also implies irreducible over Q.
bool hasRootInUpperHalfPlane(ivec3 abc) {
    int disc = abc.y*abc.y - 4*abc.x*abc.z;
    return (abc.x != 0) && (disc < 0);
}

bool march(vec3 dir, float radius) {
    float speed = length(dir);
    float t = T_MIN;
    for (int i = 0; i < MAX_STEPS; i++) {
        vec3 p = t * dir;
        float d = sdfLatticeBalls(p, radius);
        if (d < HIT_THRESHOLD) {
            ivec3 abc = ivec3(round(p));
            if (hasRootInUpperHalfPlane(abc))
                return true;
            // Invalid point: step past it
            t += (radius + HIT_THRESHOLD) / speed;
        } else {
            t += d / speed;
        }
        if (t > T_MAX) break;
    }
    return false;
}
\end{lstlisting}

For each pixel $z$, we compute $\mathbf{d} = (1, -2s, n)$ and call \texttt{march}.
A hit means $z$ is near a quadratic irrational; color it black.

We do need to check that the lattice point we hit actually gives a quadratic with roots in the upper half plane ($a \neq 0$, discriminant $< 0$).
If not, we step past and keep marching.
The discriminant check does double duty: it validates that the roots are complex, and it gives us irreducibility over $\Q$ for free, since a quadratic with negative discriminant can't factor.

\begin{figure}[ht]
    \centering
    \includegraphics[width=0.9\textwidth]{img/quadratics.png}
    \caption{Roots of integer quadratics, rendered with \texttt{BALL\_RADIUS = 0.04}. }
    \label{fig:quadratic}
\end{figure}


%% ================================================================
\section{Cubics}
\label{sec:cubic}
%% ================================================================

\subsection{The geometry of cubic coefficient space}

A cubic $ax^3 + bx^2 + cx + d$ with root $z$ factors as
\[
    a \cdot (x^2 - 2sx + n)(x - r),
\]
where $r \in \R$ is the remaining real root.
Now we have two free parameters: the leading coefficient $a$ and the real root $r$.
In $(a, b, c, d)$-space this traces out a plane through the origin:
\[
    \mathbf{P}(a, r) = a\bigl(-2s,\; n,\; 0\bigr) + ar \cdot \bigl(-1,\; 2s,\; -n\bigr),
\]
where we've written only the lower three coefficients $(b, c, d)$ since $a$ itself is the first coordinate.

So we've gone from tracing a ray through $\Z^3$ (quadratics) to finding integer points near a plane in $\Z^4$ (cubics).
Tracing a plane directly is harder.
We need to reduce.

\subsection{Monic cubics: back to a line}

Simplest case: fix $a = 1$.
Now the only free parameter is the real root $r$, and the coefficients trace an affine line in $(b, c, d)$-space:
\[
    \mathbf{L}(r) = (-2s,\; n,\; 0) + r \cdot (-1,\; 2s,\; -n).
\]
This is exactly the same setup as the quadratic case --- a one-dimensional trace through a three-dimensional integer lattice --- except the line does not pass through the origin and $r$ ranges over all of $\R$ rather than just $r > 0$.
We march in both directions from the base point.

The trace parameter $r$ has a direct meaning: it is the real root of the cubic.
This will be crucial for the irreducibility check.

\subsection{Irreducibility}

A cubic over $\Q$ is reducible if and only if it has a rational root (any factorization must include a linear factor).
We already know the real root: it is our trace parameter $r$.
So for a monic cubic, reducibility is equivalent to $r \in \Z$ --- we just check whether $r$ is near an integer.

More carefully: at a hit, we use $r$ to identify the candidate integer root $m = \lfloor r \rceil$, then verify exactly against the integer polynomial we found:
\[
    m^3 + bm^2 + cm + d = 0 \quad \Longrightarrow \quad \text{reducible; skip.}
\]
If the evaluation is nonzero, the cubic is irreducible and we keep the hit.
This check is essentially free.

\begin{figure}[ht]
    \centering
    \includegraphics[width=0.9\textwidth]{img/monic-cubics.png}
    \caption{Roots of irreducible monic cubics, after filtering out reducible cubics whose real root is an integer (which would make $z$ merely a quadratic irrational).}
    \label{fig:cubic-monic}
\end{figure}

\subsection{General cubics: slicing the plane into lines}

For general cubics, we need the full plane parameterized by $(a, r)$.
The key observation: if we fix $a$ to be a specific positive integer, the remaining family is again a line parameterized by $r$:
\[
    \mathbf{L}_a(r) = a(-2s,\; n,\; 0) + r \cdot a(-1,\; 2s,\; -n).
\]
This is a line in $(b, c, d)$-space --- the same structure as the monic case, just scaled by $a$.
(We only need positive $a$: negating $a$ negates all coefficients without changing the roots.)

So fixing $a \in \{1, 2, 3, \ldots, A_{\max}\}$ slices the plane into a family of parallel lines (parallel because the direction $(-1, 2s, -n)$ is the same up to scale), and we trace each one.

But does this actually find all the integer points near the plane?
Yes: any integer point $(a, b, c, d) \in \Z^4$ near the plane has, in particular, an integer first coordinate $a$.
So it lies near the line $\mathbf{L}_a$ for that specific value of $a$.
As long as we search over enough values of $a$, we find every nearby lattice point.

\begin{lstlisting}
// Only positive a: negating a negates all coefficients,
// which doesn't change the roots of the polynomial.
for (int a = 1; a <= A_MAX; a++) {
    vec3 base = float(a) * baseUnit;
    vec3 dir  = float(a) * dirUnit;
    if (traceLine(base, dir, R_MAX, BALL_RADIUS, a)) {
        // Darker = smaller a (simpler polynomial)
        float brightness = float(a - 1) / float(A_MAX);
        col = vec3(brightness);
        break;
    }
}
\end{lstlisting}

Setting \texttt{A\_MAX = 1} recovers the monic case.
We search from smallest $a$ outward and color by $a$: darker for simpler polynomials.

For the irreducibility check, we generalize slightly.
By the rational root theorem, any rational root of $ax^3 + bx^2 + cx + d$ has the form $p/q$ with $q \mid a$.
We already know $r$, so for each divisor $q$ of $a$ we round $qr$ to the nearest integer $p$ and evaluate the polynomial at $p/q$ in homogeneous form (multiplying through by $q^3$ to stay in integers).
Float arithmetic avoids overflow, and a threshold of $0.5$ is safe: any nonzero value of $aq^3 \cdot f(p/q)$ is at least $1$ in absolute value, so values below $0.5$ must be zero.

\begin{lstlisting}
bool isReducible(int a, float r, ivec3 bcd) {
    int b = bcd.x, c = bcd.y, d = bcd.z;
    int absA = abs(a);
    for (int q = 1; q <= absA; q++) {
        if (absA % q != 0) continue;
        float fq = float(q);
        float fn = round(fq * r);
        // Float arithmetic to avoid int overflow.
        float fval = float(a)*fn*fn*fn + float(b)*fn*fn*fq
                   + float(c)*fn*fq*fq + float(d)*fq*fq*fq;
        if (abs(fval) < 0.5) return true;
    }
    return false;
}
\end{lstlisting}

With $a \leq 5$ this runs at most a handful of iterations.

\begin{figure}[ht]
    \centering
  \includegraphics[width=0.9\textwidth]{img/cubic-all.png}
    \caption{Irreducible cubic roots with $a \leq 5$.}
    \label{fig:cubic-general}
\end{figure}


%% ================================================================
\section{Quartics}
\label{sec:quartic}
%% ================================================================

\subsection{Three dimensions of freedom}

A quartic $ax^4 + Bx^3 + Cx^2 + Dx + E$ with root $z$ factors as
\[
    a \cdot (x^2 - 2sx + n)(x^2 + px + q),
\]
with three free real parameters: the leading coefficient $a$ and the coefficients $p, q$ of the second quadratic factor.
(Whether the second factor has real or complex roots depends on the sign of $p^2 - 4q$; both cases are parameterized simultaneously.)
The coefficients are affine in $(a, p, q)$:
\begin{align*}
    B &= a(p - 2s), \\
    C &= a(q - 2sp + n), \\
    D &= a(np - 2sq), \\
    E &= anq.
\end{align*}
For fixed $a$, this traces a 2-dimensional plane in $(B, C, D, E)$-space:
\[
    \mathbf{P}_a(p,q) = a(-2s,\; n,\; 0,\; 0)
    + p \cdot a(1,\; -2s,\; n,\; 0)
    + q \cdot a(0,\; 1,\; -2s,\; n).
\]
Note the shift pattern: the two basis vectors are both $a(1, -2s, n)$, displaced by one coordinate slot.

So for fixed $a$, we need integer points in $\Z^4$ near a 2-dimensional plane.
Same trick: fix another coordinate to be an integer, reducing the plane to lines.

\subsection{Fixing two coefficients}

The first coordinate $B = a(p - 2s)$ determines $p$ once $a$ and $B$ are fixed:
\[
    p = \frac{B}{a} + 2s.
\]
Substituting, the remaining three coordinates $(C, D, E)$ become linear in $q$:
\[
    \begin{pmatrix} C \\ D \\ E \end{pmatrix}
    = a\begin{pmatrix} n - 2sp \\ np \\ 0 \end{pmatrix}
    + q \cdot a\begin{pmatrix} 1 \\ -2s \\ n \end{pmatrix}.
\]
This is a line in $\R^3$ parameterized by $q$ --- exactly the cubic setup.
The direction $a(1, -2s, n)$ is the same for every $B$; only the base point shifts.

Same completeness argument as before: any integer point $(B, C, D, E) \in \Z^4$ near the plane has an integer $B$, so it lies near one of our lines.
The algorithm is a double loop:

\begin{lstlisting}
// Only positive a: negating all coefficients
// doesn't change the roots.
for (int a = 1; a <= A_MAX; a++) {
    float fa = float(a);
    vec3 dirA = fa * vec3(1.0, -2.0 * s, n);
    for (int B = -B_RANGE; B <= B_RANGE; B++) {
        float p = (float(B) + 2.0 * fa * s) / fa;
        vec3 base = fa * vec3(n - 2.0*s*p, n*p, 0.0);
        result = traceLine(base, dirA, Q_MAX,
                           BALL_RADIUS, a, B, p, s, n);
        if (result != NO_HIT) break;
    }
    if (result != NO_HIT) break;
}
\end{lstlisting}

Setting \texttt{A\_MAX = 1} restricts to monic quartics.
The cost per pixel is $A_{\max} \times (2 B_{\text{range}} + 1)$ traces, but the early exit on first hit means most pixels terminate much sooner.

\subsection{Irreducibility}

Quartic irreducibility is trickier than cubics, because there are two ways a quartic can be reducible: it can have a rational root, or it can factor into two integer quadratics.

For rational roots: if $E = 0$ then $x = 0$ is a root and we're done.
Otherwise, any rational root has to come from the second factor $x^2 + px + q$ (the first factor's roots are $z$ and $\bar{z}$, which aren't real).
If $p^2 - 4q < 0$ the second factor has no real roots either, so there's nothing to check.
Otherwise we get two real root candidates from the quadratic formula and check whether either is rational with denominator dividing $a$, same strategy as the cubic case: for each divisor $d$ of $a$, round $d \cdot r_i$ to the nearest integer and evaluate the quartic in homogeneous form ($d^4$ times the polynomial at $r_i$) to avoid overflow.
Threshold $0.5$ is safe since a nonzero evaluation has magnitude at least $1/d^4$.
With $a \leq 3$, this checks at most 4 candidates.

For factorization into two integer quadratics: we're looking for
\[
    (a_1 x^2 + Px + Q)(a_2 x^2 + Rx + S), \qquad a_1 a_2 = a, \quad P, Q, R, S \in \Z.
\]
Brute-forcing over $P$ is expensive and needs an arbitrary range cutoff.
But we know $z$, so any integer quadratic factor containing $z$ and $\bar{z}$ must have $P \approx -2a_1 s$ and $Q \approx a_1 n$.
For each factorization $a = a_1 a_2$, we try $P$ in a window of three around $\lfloor -2a_1 s \rceil$.
I initially just rounded to the nearest integer, but that misses factorizations when $-2a_1 s$ is near a half-integer --- you get false-positive blue triangles around the quadratic irrationals.
Trying $\pm 1$ around the estimate fixes it.

For each candidate $P$, the rest is determined algebraically.
$R = (B - a_2 P)/a_1$ (requiring divisibility).
Setting $\sigma = C - PR$, the values $a_2 Q$ and $a_1 S$ are roots of $t^2 - \sigma t + a_1 a_2 E = 0$, so we need $\sigma^2 - 4aE$ to be a perfect square.
Then $Q = (\sigma \pm \sqrt{\sigma^2 - 4aE\,})/(2a_2)$ and $S = (\sigma - a_2 Q)/a_1$, both required to be integers, and we verify $PS + QR = D$.
Iterating $a_1$ from $1$ to $|a|$ automatically covers both assignments of which factor contains $z$.

\subsection{Root-type coloring}

Once a quartic passes both checks, we can classify it by root type.
The first factor $x^2 - 2sx + n$ always has complex roots (its discriminant is $4s^2 - 4n = -4y^2 < 0$ since $z$ is in the upper half plane).
The second factor $x^2 + px + q$ may have either real or complex roots, depending on the sign of $p^2 - 4q$.

When $p^2 - 4q < 0$, \emph{both} quadratic factors have complex roots, so all four roots of the quartic are non-real (two conjugate pairs).
When $p^2 - 4q \geq 0$, the second factor contributes two real roots, giving two complex and two real roots in total.

We color the all-complex case \textcolor{blue}{\textbf{blue}} and the mixed case black, giving a visual separation of root type.
For example, $x^4 + 1$ factors over $\R$ as $(x^2 + \sqrt{2}\,x + 1)(x^2 - \sqrt{2}\,x + 1)$, where both factors have negative discriminant; this quartic would be colored blue.

\begin{figure}[ht]
    \centering
    % \includegraphics[width=0.7\textwidth]{quartic-monic.png}
\includegraphics[width=0.9\textwidth]{img/monic-quartic.png}
    \caption{Irreducible monic quartics. Black: two complex and two real roots. Blue: four complex roots (both quadratic factors have complex roots).}
    \label{fig:quartic-monic}
\end{figure}

\begin{figure}[ht]
    \centering
\includegraphics[width=0.9\textwidth]{img/quartics.png}
    \caption{Irreducible quartics with $a \leq 3$, with the blue/black root-type distinction preserved and brightness indicating leading coefficient.}
    \label{fig:quartic-general}
\end{figure}


%% ================================================================
\section{Summary and the General Pattern}
\label{sec:summary}
%% ================================================================

The whole thing comes down to two observations.
The family of degree-$d$ polynomials with a prescribed complex root $z$ forms a $(d-1)$-dimensional linear subspace in coefficient space.
And we can reduce finding integer points near a $k$-dimensional subspace to a family of one-dimensional raymarches by fixing $k - 1$ coordinates to be integers.

\begin{center}
\begin{tabular}{ccclp{4.5cm}}
    \toprule
    Degree & Subspace dim.\ & Coords.\ fixed & Traces per pixel & Irreducibility check \\
    \midrule
    2 & 1 & 0 & 1 & None needed (disc $< 0$ suffices) \\
    3 & 2 & 1 ($= a$) & $A_{\max}$ & Rational root via known $r$ \\
    4 & 3 & 2 ($= a, B$) & $A_{\max} \times (2B_{\text{range}} + 1)$ & Rational root $+$ integer quadratic factorization \\
    \bottomrule
\end{tabular}
\end{center}

This pattern continues: degree $d$ means fixing $d - 2$ integer coordinates and tracing a $(d-2)$-fold family of lines.
At each hit, the known factorization gives us cheap irreducibility checks: for cubics we know the real root and just check if it's rational; for quartics we know the second quadratic factor and can test for rational roots and integer factorizations.
For quartics we also color by root type --- blue for four complex roots, black for two complex and two real --- using the sign of the second factor's discriminant.


\appendix

%% ================================================================
\section{Quadratic Shader}
\label{app:quadratic}
%% ================================================================

\begin{lstlisting}
// ============================================================
// QUADRATIC IRRATIONALS IN THE UPPER HALF PLANE
//
// For z in the upper half plane, the set of real quadratics
// with z as a root forms a ray in coefficient space (a,b,c):
//
//     R(t) = t * (1, -2*Re(z), |z|^2)    for t > 0
//
// We raymarch this ray looking for near-misses with integer
// points, which correspond to quadratic irrationals near z.
// ============================================================

// --- Parameters ---
const float BALL_RADIUS = 0.04;
const float T_MIN = 0.5;
const float T_MAX = 1000.0;
const float HIT_THRESHOLD = 0.001;
const int   MAX_STEPS = 300;

// --- View settings ---
const vec2 VIEW_CENTER = vec2(0.0, 1.0);
const float VIEW_HEIGHT = 2.0;

// ============================================================
// COEFFICIENT SPACE GEOMETRY
// ============================================================

vec3 rayDirection(vec2 z) {
    float x = z.x;
    float normSq = dot(z, z);
    return vec3(1.0, -2.0 * x, normSq);
}

float sdfLatticeBalls(vec3 p, float radius) {
    vec3 offset = p - round(p);
    return length(offset) - radius;
}

ivec3 nearestInteger(vec3 p) {
    return ivec3(round(p));
}

// Check that (a,b,c) gives a root in the upper half plane.
// Requires: a != 0 and discriminant < 0 (complex roots).
// disc < 0 also implies irreducible over Q.
bool hasRootInUpperHalfPlane(ivec3 abc) {
    int a = abc.x;
    int b = abc.y;
    int c = abc.z;
    int disc = b*b - 4*a*c;
    return (a != 0) && (disc < 0);
}

// ============================================================
// RAYMARCHING
// ============================================================

bool march(vec3 dir, float radius, out ivec3 hitCoeffs) {
    float speed = length(dir);
    float t = T_MIN;

    for (int i = 0; i < MAX_STEPS; i++) {
        vec3 p = t * dir;
        float d = sdfLatticeBalls(p, radius);

        if (d < HIT_THRESHOLD) {
            ivec3 abc = nearestInteger(p);
            if (hasRootInUpperHalfPlane(abc)) {
                hitCoeffs = abc;
                return true;
            }
            // Invalid point: step past it
            t += (radius + HIT_THRESHOLD) / speed;
        } else {
            t += d / speed;
        }

        if (t > T_MAX) break;
    }

    return false;
}

// ============================================================
// VISUALIZATION
// ============================================================

vec2 pixelToHalfPlane(vec2 fragCoord, vec2 resolution) {
    vec2 uv = fragCoord / resolution;
    float aspect = resolution.x / resolution.y;
    vec2 z;
    z.x = VIEW_CENTER.x
        + (uv.x - 0.5) * VIEW_HEIGHT * aspect;
    z.y = VIEW_CENTER.y
        + (uv.y - 0.5) * VIEW_HEIGHT;
    return z;
}

void mainImage(out vec4 fragColor, in vec2 fragCoord) {
    vec2 z = pixelToHalfPlane(fragCoord, iResolution.xy);

    if (z.y <= 0.0) {
        fragColor = vec4(vec3(0.3), 1.0);
        return;
    }

    vec3 dir = rayDirection(z);

    ivec3 coeffs;
    bool hit = march(dir, BALL_RADIUS, coeffs);

    vec3 col = hit ? vec3(0.0) : vec3(1.0);
    fragColor = vec4(col, 1.0);
}
\end{lstlisting}


%% ================================================================
\section{Cubic Shader}
\label{app:cubic}
%% ================================================================

\begin{lstlisting}
// ============================================================
// CUBIC ALGEBRAIC NUMBERS IN THE UPPER HALF PLANE
//
// A cubic with leading coefficient a and root z factors as
//
//   a * (x^2 - 2sx + n)(x - r)
//
// where s = Re(z), n = |z|^2. Fixing integer a, the lower
// three coefficients (b,c,d) trace a line parameterized by
// the real root r:
//
//   L_a(r) = a(-2s, n, 0) + r * a(-1, 2s, -n)
//
// We loop over a = 1, ..., A_MAX and raymarch each line in
// both directions. Only positive a: negating a negates all
// coefficients without changing the roots.
//
// Set A_MAX = 1 for monic cubics only.
// ============================================================

// --- Search parameters ---
const float BALL_RADIUS = 0.02;
const float R_MAX = 50.0;
const float HIT_THRESHOLD = 0.001;
const int   MAX_STEPS = 256;
const int   A_MAX = 5;

// --- View ---
const vec2 VIEW_CENTER = vec2(0.0, 1.5);
const float VIEW_HEIGHT = 3.0;

const bool IRREDUCIBLE_ONLY = true;

// ============================================================
// Lattice SDF
// ============================================================

float sdfLatticeBalls(vec3 p, float radius) {
    return length(p - round(p)) - radius;
}

// ============================================================
// Irreducibility
//
// The trace parameter r IS the real root of the cubic.
// A cubic over Q is reducible iff it has a rational root,
// and by the rational root theorem any such root has the
// form p/q with q | a.
//
// So for each divisor q of a, we round q*r to the nearest
// integer and evaluate the polynomial there. If it vanishes,
// the cubic is reducible. Uses float arithmetic to avoid
// int overflow; threshold 0.5 is safe because any nonzero
// value of an integer polynomial at a rational point is
// at least 1/q^3 in absolute value.
// ============================================================

bool isReducible(int a, float r, ivec3 bcd) {
    int b = bcd.x, c = bcd.y, d = bcd.z;
    int absA = abs(a);

    for (int q = 1; q <= absA; q++) {
        if (absA % q != 0) continue;

        float fq = float(q);
        float fn = round(fq * r);

        float fval = float(a) * fn * fn * fn
                   + float(b) * fn * fn * fq
                   + float(c) * fn * fq * fq
                   + float(d) * fq * fq * fq;

        if (abs(fval) < 0.5) return true;
    }

    return false;
}

// ============================================================
// Line tracing
//
// March along base + t*dir for t >= 0.
// The trace parameter t = |r|; the sign argument recovers
// the actual real root r = sign * t.
// ============================================================

bool traceHalf(vec3 base, vec3 dir, float tMax,
               float radius, int a, float sign) {
    float speed = length(dir);
    float t = 0.0;

    for (int i = 0; i < MAX_STEPS; i++) {
        vec3 p = base + t * dir;
        float d = sdfLatticeBalls(p, radius);

        if (d < HIT_THRESHOLD) {
            if (IRREDUCIBLE_ONLY) {
                float r = sign * t;
                ivec3 bcd = ivec3(round(p));
                if (isReducible(a, r, bcd)) {
                    t += (radius + HIT_THRESHOLD)
                       / speed;
                    continue;
                }
            }
            return true;
        }

        t += d / speed;
        if (t > tMax) break;
    }

    return false;
}

bool traceLine(vec3 base, vec3 dir, float tMax,
               float radius, int a) {
    if (traceHalf(base,  dir, tMax, radius, a,  1.0))
        return true;
    if (traceHalf(base, -dir, tMax, radius, a, -1.0))
        return true;
    return false;
}

// ============================================================
// Coordinate map
// ============================================================

vec2 pixelToHalfPlane(vec2 fragCoord, vec2 resolution) {
    vec2 uv = fragCoord / resolution;
    float aspect = resolution.x / resolution.y;
    return VIEW_CENTER
         + (uv - 0.5)
         * vec2(VIEW_HEIGHT * aspect, VIEW_HEIGHT);
}

// ============================================================
// Main
// ============================================================

void mainImage(out vec4 fragColor, in vec2 fragCoord) {
    vec2 z = pixelToHalfPlane(fragCoord, iResolution.xy);

    if (z.y <= 0.0) {
        fragColor = vec4(vec3(0.5), 1.0);
        return;
    }

    float s = z.x;
    float n = dot(z, z);

    vec3 baseUnit = vec3(-2.0 * s, n, 0.0);
    vec3 dirUnit  = vec3(-1.0, 2.0 * s, -n);

    vec3 col = vec3(1.0);

    for (int a = 1; a <= A_MAX; a++) {
        vec3 base = float(a) * baseUnit;
        vec3 dir  = float(a) * dirUnit;

        if (traceLine(base, dir, R_MAX,
                      BALL_RADIUS, a)) {
            // Darker = smaller a (simpler polynomial)
            float brightness =
                float(a - 1) / float(A_MAX);
            col = vec3(brightness);
            break;
        }
    }

    fragColor = vec4(col, 1.0);
}
\end{lstlisting}


%% ================================================================
\section{Quartic Shader}
\label{app:quartic}
%% ================================================================

\begin{lstlisting}
// ============================================================
// QUARTIC ALGEBRAIC NUMBERS IN THE UPPER HALF PLANE
//
// A quartic with leading coefficient a and root z factors as
//
//   a * (x^2 - 2sx + n)(x^2 + px + q)
//
// where s = Re(z), n = |z|^2. Fix integer a, then fix
// integer B = a(p - 2s), giving p = B/a + 2s. The remaining
// three coefficients (C,D,E) trace a line parameterized by q:
//
//   L(q) = a(n - 2sp, np, 0) + q * a(1, -2s, n)
//
// Set A_MAX = 1 for monic only.
//
// Colors:
//   BLACK/GRAY - irreducible, 2 complex + 2 real roots
//   BLUE       - irreducible, 4 complex roots
//   WHITE      - no hit
// ============================================================

// --- Search parameters ---
const float BALL_RADIUS = 0.005;
const float Q_MAX = 50.0;
const float HIT_THRESHOLD = 0.001;
const int   MAX_STEPS = 128;

const int A_MAX = 3;
const int B_RANGE = 10;

// --- View ---
const vec2 VIEW_CENTER = vec2(0.0, 1.5);
const float VIEW_HEIGHT = 3.0;

const bool IRREDUCIBLE_ONLY = true;

const int NO_HIT = 0;
const int HIT_GENERIC = 1;        // 2 complex + 2 real
const int HIT_ALL_COMPLEX = 2;    // 4 complex roots

// ============================================================
// Lattice SDF
// ============================================================

float sdfLatticeBalls(vec3 pt, float radius) {
    return length(pt - round(pt)) - radius;
}

// ============================================================
// Irreducibility checks
//
// A quartic over Q is reducible iff:
//   1. It has a rational root (linear factor), or
//   2. It factors as two quadratics over Z.
// There is no other case.
// ============================================================

// Evaluate ax^4+Bx^3+Cx^2+Dx+E at x = fn/fd in
// homogeneous form (multiply through by fd^4).
// Float arithmetic to avoid int overflow.
float evalQuarticHom(int a, int B, int C, int D, int E,
                     float fn, float fd) {
    return float(a) * fn * fn * fn * fn
         + float(B) * fn * fn * fn * fd
         + float(C) * fn * fn * fd * fd
         + float(D) * fn * fd * fd * fd
         + float(E) * fd * fd * fd * fd;
}

// --- Check 1: rational linear factor ---
//
// If p^2 - 4q < 0, the second factor has no real roots,
// so no linear factor exists. Otherwise, check whether
// either real root is rational with denominator | a.
// Threshold 0.5 is safe: nonzero evaluations of integer
// polynomials at rationals have magnitude >= 1/denom^4.

bool hasRationalRoot(int a, int B, int C, int D, int E,
                     float p, float q) {
    if (E == 0) return true;

    float disc = p * p - 4.0 * q;
    if (disc < 0.0) return false;

    float sqrtDisc = sqrt(disc);
    float r1 = (-p + sqrtDisc) / 2.0;
    float r2 = (-p - sqrtDisc) / 2.0;

    int absA = abs(a);
    for (int div = 1; div <= absA; div++) {
        if (absA % div != 0) continue;
        float fd = float(div);

        float fn = round(fd * r1);
        if (abs(evalQuarticHom(a,B,C,D,E,fn,fd)) < 0.5)
            return true;

        fn = round(fd * r2);
        if (abs(evalQuarticHom(a,B,C,D,E,fn,fd)) < 0.5)
            return true;
    }
    return false;
}

// --- Check 2: product of two integer quadratics ---
//
// We want: (a1*x^2+P*x+Q)(a2*x^2+R*x+S) with a1*a2=a.
// We know z, so any integer quadratic containing z must
// have P ~ -2*a1*s. Try a window of 3 values around
// the pixel-based estimate to handle rounding at
// half-integers. The rest is determined algebraically.
//
// Looping a1 from 1 to |a| automatically covers both
// assignments of which factor contains z.

bool factorsAsIntegerQuadratics(
        int a, int B, int C, int D, int E,
        float s, float n) {
    int absA = abs(a);

    for (int a1 = 1; a1 <= absA; a1++) {
        if (absA % a1 != 0) continue;
        int a2 = a / a1;

        // Pixel-based estimate for P
        float fP = -2.0 * float(a1) * s;
        int Pest = int(round(fP));

        // Try Pest-1, Pest, Pest+1
        for (int dP = -1; dP <= 1; dP++) {
            int P = Pest + dP;

            // x^3 coeff: a1*R + a2*P = B
            int Rnum = B - a2 * P;
            if (Rnum % a1 != 0) continue;
            int R = Rnum / a1;

            // sigma = C - P*R; then a2*Q and a1*S
            // are roots of t^2 - sigma*t + a*E = 0
            int sigma = C - P * R;
            int disc = sigma*sigma - 4*a*E;
            if (disc < 0) continue;

            int sqrtDisc =
                int(round(sqrt(float(disc))));
            if (sqrtDisc*sqrtDisc != disc) continue;

            for (int sgn = -1; sgn <= 1; sgn += 2) {
                int Qnum = sigma + sgn * sqrtDisc;
                if (Qnum % (2*a2) != 0) continue;
                int Q = Qnum / (2 * a2);

                int Snum = sigma - a2 * Q;
                if (Snum % a1 != 0) continue;
                int S = Snum / a1;

                // Verify x^1 coefficient
                if (P*S + Q*R == D) return true;
            }
        }
    }
    return false;
}

bool isReducible(int a, int B, int C, int D, int E,
                 float p, float q, float s, float n) {
    if (hasRationalRoot(a, B, C, D, E, p, q))
        return true;
    if (factorsAsIntegerQuadratics(a, B, C, D, E, s, n))
        return true;
    return false;
}

// ============================================================
// Line tracing
//
// March along base + t*dir for t >= 0.
// Trace parameter t = |q|; sign recovers q = sign * t.
// ============================================================

int traceHalf(vec3 base, vec3 dir, float tMax,
              float radius, int a, int B, float p,
              float sign, float s, float n) {
    float speed = length(dir);
    float t = 0.0;

    for (int i = 0; i < MAX_STEPS; i++) {
        vec3 pt = base + t * dir;
        float d = sdfLatticeBalls(pt, radius);

        if (d < HIT_THRESHOLD) {
            float q = sign * t;

            if (IRREDUCIBLE_ONLY) {
                ivec3 CDE = ivec3(round(pt));
                if (isReducible(a, B,
                        CDE.x, CDE.y, CDE.z,
                        p, q, s, n)) {
                    t += (radius + HIT_THRESHOLD)
                       / speed;
                    continue;
                }
            }

            // Root type classification:
            // First factor always has complex roots
            // (disc = 4s^2 - 4n = -4y^2 < 0).
            // Second factor x^2+px+q has complex roots
            // iff p^2-4q < 0.
            if (p*p - 4.0*q < 0.0)
                return HIT_ALL_COMPLEX;
            return HIT_GENERIC;
        }

        t += d / speed;
        if (t > tMax) break;
    }

    return NO_HIT;
}

int traceLine(vec3 base, vec3 dir, float tMax,
              float radius, int a, int B, float p,
              float s, float n) {
    int result = traceHalf(base, dir, tMax, radius,
                           a, B, p, 1.0, s, n);
    if (result != NO_HIT) return result;
    return traceHalf(base, -dir, tMax, radius,
                     a, B, p, -1.0, s, n);
}

// ============================================================
// Coordinate map
// ============================================================

vec2 pixelToHalfPlane(vec2 fragCoord,
                      vec2 resolution) {
    vec2 uv = fragCoord / resolution;
    float aspect = resolution.x / resolution.y;
    return VIEW_CENTER
         + (uv - 0.5)
         * vec2(VIEW_HEIGHT * aspect, VIEW_HEIGHT);
}

// ============================================================
// Main
// ============================================================

void mainImage(out vec4 fragColor,
               in vec2 fragCoord) {
    vec2 z = pixelToHalfPlane(fragCoord,
                              iResolution.xy);

    if (z.y <= 0.0) {
        fragColor = vec4(vec3(0.5), 1.0);
        return;
    }

    float s = z.x;
    float n = dot(z, z);

    int result = NO_HIT;
    int hitA = 0;

    for (int a = 1; a <= A_MAX; a++) {
        float fa = float(a);
        vec3 dirA = fa * vec3(1.0, -2.0 * s, n);

        for (int B = -B_RANGE; B <= B_RANGE; B++) {
            float p = (float(B)
                     + 2.0 * fa * s) / fa;
            vec3 base = fa * vec3(
                n - 2.0 * s * p, n * p, 0.0);

            result = traceLine(base, dirA, Q_MAX,
                BALL_RADIUS, a, B, p, s, n);
            if (result != NO_HIT) {
                hitA = a; break;
            }
        }
        if (result != NO_HIT) break;
    }

    vec3 col = vec3(1.0);
    if (result != NO_HIT) {
        float brightness =
            float(hitA - 1) / float(A_MAX);
        if (result == HIT_ALL_COMPLEX) {
            col = mix(vec3(0.2, 0.4, 1.0),
                      vec3(1.0), brightness);
        } else {
            col = vec3(brightness);
        }
    }

    fragColor = vec4(col, 1.0);
}
\end{lstlisting}


\end{document}